\section{Problem setting}
\label{sec:setting}

\paragraph{Equations of motion}
A rigid body $b$ has centre-of-mass position $r_b\in\mathbb R^3$, rotation $Q_b\in SO(3)$,
translational velocity $v_b=\dot r_b$, body-frame angular velocity $\omega_b$ with
$\dot Q_b=Q_b\,\widehat{\omega_b}$, mass $m_b$ and body-frame inertia $J_b$. With
$q=(r_b,Q_b)_b$, $v=(v_b,\omega_b)_b$ and $a=\dot v$, a mechanism with lower-pair joints is
governed by
\begin{align}
  M\dot v + G(v) + \Phi_q(q)^T\lambda &= F(q,v,t), \label{eq:dae-dyn}\\
  \dot q &= \mathcal K(q)v, \label{eq:dae-kin}\\
  \Phi(q) &= 0, \label{eq:dae-constraint}
\end{align}
where $M=\operatorname{diag}(m_bI,J_b)$, $G$ collects the gyroscopic terms
$\omega_b\times J_b\omega_b$, $\Phi$ stacks the joint constraints, $\lambda$ are the Lagrange
multipliers, and $\mathcal K$ is the kinematic map, which is the identity on the translational part
and $\omega_b\mapsto Q_b\widehat{\omega_b}$ on the rotational part. The system is an index-3
differential-algebraic equation on $(\mathbb R^3\times SO(3))^{N}$. Differentiating the constraint
once and twice gives the velocity- and acceleration-level constraints
\begin{equation}
  \Phi_q(q)v=0,\qquad \Phi_q(q)a+\dot\Phi_q(q,v)v=0,
  \label{eq:fullva}
\end{equation}
which the exact solution satisfies identically; the method of Section~\ref{sec:method} enforces
Eq.~\eqref{eq:dae-constraint} and both equations of Eq.~\eqref{eq:fullva} at every stage, which
is what the acronym FullVA (full velocity- and acceleration-level enforcement) refers to.

\paragraph{Rotations}
Rotations are represented by unit quaternions in the implementation and advanced by the retraction
\begin{equation}
  Q_{n+1}=Q_n\,\mathrm{Exp}(\eta_{n+1}),\qquad Q_i=Q_n\,\mathrm{Exp}(\eta_i),
  \label{eq:retraction}
\end{equation}
with $\eta\in\mathbb R^3\cong\mathfrak{so}(3)$ the local rotation increment and $\mathrm{Exp}$
the exponential map. The right-trivialized differential of the exponential map is written
$J_r(\eta)$, so that $\dot\eta=J_r^{-1}(\eta)\,\omega$ along a motion with body angular velocity
$\omega$.

\paragraph{Lower pairs and joint coordinates}
A lower pair $j$ connects a parent body $p(j)$ (possibly the ground) and a child body $c(j)$. Let
$P_j(q)$ and $C_j(q)$ be the attachment points on the two bodies, $d_j(q)=C_j(q)-P_j(q)$ the
relative displacement, $a_j(q)$ the joint axis carried by the parent, $\tilde a_j(q)$ the axis
carried by the child, and $\omega^{\rm rel}_j=Q_{c(j)}\omega_{c(j)}-Q_{p(j)}\omega_{p(j)}$ the
relative angular velocity in the world frame. Each pair carries a fixed world direction $n_j$, the
joint axis at $t=0$, and a fixed orthonormal pair $e_{j,1},e_{j,2}\perp n_j$. A pair with two
degrees of freedom (sliding along and rotating about its axis) has the joint coordinates
\begin{equation}
  s_j(q)=n_j\cdot d_j(q),\qquad
  \theta_j(q)=\operatorname{atan2}\bigl(\tilde t_j\cdot(a_j\times t_j),\ \tilde t_j\cdot t_j\bigr),
  \label{eq:joint-coordinates}
\end{equation}
where $t_j,\tilde t_j$ are twist reference vectors fixed in the parent and in the child. Their
rates are $\dot s_j=n_j\cdot\dot d_j$ and $\dot\theta_j=a_j\cdot\omega^{\rm rel}_j$, and their
second rates are $\ddot s_j=n_j\cdot\ddot d_j$ and
$\ddot\theta_j=a_j\cdot\dot\omega^{\rm rel}_j+\dot a_j\cdot\omega^{\rm rel}_j$. Revolute,
prismatic and spherical joints are obtained by fixing one or both coordinates or by dropping the
alignment rows; the constraint rows of the four benchmark mechanisms are listed in
\ref{app:lower-pair}.

\paragraph{The two-body chain}
The main benchmark is a two-body open chain (Table~\ref{tab:chain}). Body $0$ is connected to the
ground by a pair with sliding and rotation along a fixed ground axis; body $1$ is connected to body
$0$ by a second pair whose axis $a_1(q)=Q_0\hat a^{\rm next}_0$ is fixed in body $0$. The two axes
fixed in body $0$ are not parallel and body $0$ spins, so $a_1(q)$ moves on a cone about the ground
axis while $n_1$ and $e_{1,m}$ stay frozen at their $t=0$ values. The second pair is therefore a
four-constraint lower pair with a fixed sliding direction and a moving rotation axis, not a
cylindrical pair in the strict sense; its friction sliding speed is $\dot s_1\,(n_1\cdot a_1)$, and
the modelling is regular as long as $n_1\cdot a_1(q)\neq0$. On the initial state of
Table~\ref{tab:chain} this factor decreases from $1$ to $0.34$ on $[0,0.5]$ and vanishes at
$t\approx0.605$, where the frozen sliding direction ceases to define the pair; all chain
experiments are run on the regular branch $[0,0.5]$. Each joint carries gravity, constant external
loads, viscous damping and a regularized Brown--McPhee friction torque and force
\cite{browmcphee2016friction} with static and dynamic coefficients $\mu_s,\mu_d$ and Stribeck
velocity $v_s$; the friction force enters the Newton--Euler rows through the multipliers (normal
force) and the sliding speed, and is $C^\infty$ in the state for $v_s>0$. The smooth case uses
$v_s=0.5$; Section~\ref{sec:numerical} varies $v_s$ down to $0.02$.

\begin{table}[htbp]
\centering
\caption{Two-body chain benchmark. Directions are given before normalization; the initial rotations
are the unique rotations mapping the body-fixed axis and twist vectors of each pair onto their
parent-side counterparts.}
\label{tab:chain}
\small
\begin{tabular}{P{0.30\linewidth}P{0.60\linewidth}}
\toprule
Quantity & Value \\
\midrule
Masses $m_0,m_1$ [kg] & $2.7$, $1.9$ \\
Inertias $J_0,J_1$ [kg\,m$^2$] & $\operatorname{diag}(0.11,0.24,0.31)$, $\operatorname{diag}(0.08,0.19,0.27)$ \\
Ground axis, ground twist & $(0.31,0.78,0.54)$, first vector of an orthonormal complement \\
Body-fixed axes (parent side) $\hat a^{\rm prev}_0,\hat a^{\rm prev}_1$ & $(0.18,0.91,-0.37)$, $(-0.24,0.87,0.43)$ \\
Body-fixed axes (child side) $\hat a^{\rm next}_0,\hat a^{\rm next}_1$ & $(0.42,0.69,0.59)$, $(0,1,0)$ \\
Attachment points $s^{\rm prev}_0,s^{\rm prev}_1$ [m] & $(0.12,-0.04,0.18)$, $(-0.08,0.05,0.16)$ \\
Attachment points $s^{\rm next}_0,s^{\rm next}_1$ [m] & $(0.16,0.02,-0.20)$, $(0,0,0)$ \\
Gravity [m/s$^2$] & $(0,0,-9.81)$ \\
External forces (world) [N] & $(0.8,-0.35,0.25)$, $(-0.25,0.45,-0.10)$ \\
External torques (body) [N\,m] & $(0.03,-0.02,0.015)$, $(-0.018,0.026,-0.014)$ \\
Friction $\mu_s,\mu_d$, viscous coefficient, radius & $0.30$, $0.20$, $0.030$, $1.0$ \\
Stribeck velocity $v_s$ & $0.5$ (smooth case); $0.2,0.1,0.05,0.02$ in the sweep \\
Initial slide $s_j(0)$, slide rate $\dot s_j(0)$, spin rate $\dot\theta_j(0)$ & $(0.35,0.42)$, $(0.42,-0.26)$, $(0.18,-0.14)$ \\
\bottomrule
\end{tabular}
\end{table}

\paperfig{\figpath/chain_schematic.png}
{The two-body chain: body $0$ slides and spins along the fixed ground axis, body $1$ slides and spins
along the axis $a_1(q)$ carried by body $0$; the second pair keeps its sliding direction $n_1$ and
its basis $e_{1,m}$ frozen at their $t=0$ values.}
{fig:chain}

\paragraph{Notation}
Throughout, $h$ is the step size, $t_n=nh$, $x_n=(q_n,v_n)$ the endpoint state (accelerations and
multipliers are stage quantities), $c_i,b_i,A_{ik}$ ($i,k=1,2,3$) the nodes, weights and matrix of
the three-stage Gauss tableau, $Z$ the stacked stage vector, and $\|\cdot\|$ a fixed norm on the
finite-dimensional space at hand. Constants $C$ with subscripts are independent of $h$ and of the
number of steps.

\section{The \method{} step}
\label{sec:method}

\subsection{Stage unknowns and stage system}

For each stage $i=1,2,3$ and each body $b$ the unknowns are the translational position $r_{b,i}$,
the rotation increment $\eta_{b,i}$ with $Q_{b,i}=Q_{b,n}\mathrm{Exp}(\eta_{b,i})$, the velocities
$v_{b,i},\omega_{b,i}$ and the accelerations $a_{b,i},\alpha_{b,i}$; each lower pair $j$
contributes its multipliers $\lambda_{j,i}$, four per pair for the chain. With two bodies and two
pairs the stage vector
\begin{equation}
  Z=(Z_1,Z_2,Z_3),\qquad Z_i=(r_i,\eta_i,v_i,\omega_i,a_i,\alpha_i,\lambda_i),
  \label{eq:stage-variable-block}
\end{equation}
has $3\,(2\cdot18+2\cdot4)=132$ components.

Write $\psi_j(q)=\tilde a_j(q)-a_j(q)$ for the axis-alignment residual and, for a joint-coordinate
function $g$ with rate $\dot g$,
\begin{equation}
  \Delta_i[g]=g(Z_i)-g(x_n)-h\sum_{k=1}^3A_{ik}\,\dot g(Z_k)
  \label{eq:collocation-defect-operator}
\end{equation}
for its Gauss collocation defect at stage $i$. With $t_i=t_n+c_ih$, the applied wrench
$(f_{b,i},\tau_{b,i})$ evaluated at the stage, and the total force and torque
$F^r_{b,i}=f_{b,i}+\Phi_{r,b}(q_i)^T\lambda_i$, $T^\theta_{b,i}=\tau_{b,i}+\Phi_{\theta,b}(q_i)^T\lambda_i$
including the multiplier wrench, the stage system reads, for every pair $j$, every $m=1,2$ and
every body $b$,
\begin{equation}
\label{eq:g6fullva-expanded-stage}
\begin{aligned}
0 &= e_{j,m}\cdot d_j(q_i), &
0 &= e_{j,m}\cdot\psi_j(q_i),\\
0 &= e_{j,m}\cdot\dot d_j(q_i,v_i), &
0 &= e_{j,m}\cdot\dot\psi_j(q_i,v_i),\\
0 &= e_{j,m}\cdot\ddot d_j(q_i,v_i,a_i), &
0 &= e_{j,m}\cdot\ddot\psi_j(q_i,v_i,a_i),\\
0 &= \Delta_i[s_j], & 0 &= \Delta_i[\dot s_j],\\
0 &= \Delta_i[\theta_j], & 0 &= \Delta_i[\dot\theta_j],\\
0 &= m_b a_{b,i}-F^r_{b,i}, &
0 &= J_b\alpha_{b,i}+\omega_{b,i}\times J_b\omega_{b,i}-T^\theta_{b,i} .
\end{aligned}
\end{equation}
The first three lines are the lower-pair constraints at position, velocity and acceleration level
($12$ rows per pair and stage); the fourth and fifth lines are Gauss collocation of the four joint
coordinates $(s_j,\dot s_j,\theta_j,\dot\theta_j)$ ($4$ rows per pair and stage); the last line is
the Newton--Euler balance ($6$ rows per body and stage). Table~\ref{tab:rows} gives the inventory:
$44$ rows per stage for the $44$ unknowns, $72+24+36=132$ over the three stages. The absolute body
coordinates $(r_i,\eta_i,v_i,\omega_i)$ carry no collocation row of their own; they are determined
by the constraint rows once the joint coordinates are collocated. The implementation projects the
two translational collocation rows onto the current joint axis $a_j(q_i)$ rather than onto $n_j$;
on the constraint manifold $d_j,\dot d_j,\ddot d_j\parallel n_j$, so the two forms differ by the
factor $n_j\cdot a_j(q_i)$, nonzero on the regular branch, and define the same stage solution.
Driven mechanisms add prescribed right-hand sides to the three constraint levels of the driven
row. The multiplier wrench collects the joint normal forces, the axis-alignment torques and the
friction force along the joint axis; the friction law enters only there.

\begin{table}[htbp]
\centering
\caption{Row inventory of the stage system for the two-body chain (per stage and in total).}
\label{tab:rows}
\small
\begin{tabular}{P{0.46\linewidth}P{0.16\linewidth}P{0.10\linewidth}P{0.10\linewidth}}
\toprule
Row family & Per pair or body & Per stage & Total \\
\midrule
Position-level constraints $\Phi=0$ & $4$ per pair & $8$ & $24$ \\
Velocity-level constraints $\Phi_qv=0$ & $4$ per pair & $8$ & $24$ \\
Acceleration-level constraints $\Phi_qa+\dot\Phi_qv=0$ & $4$ per pair & $8$ & $24$ \\
Joint-coordinate collocation $\Delta_i[s_j],\Delta_i[\dot s_j],\Delta_i[\theta_j],\Delta_i[\dot\theta_j]$ & $4$ per pair & $8$ & $24$ \\
Newton--Euler balance & $6$ per body & $12$ & $36$ \\
\midrule
Unknowns $(r,\eta,v,\omega,a,\alpha)$ per body, $\lambda$ per pair & $18$, $4$ & $44$ & $132$ \\
\bottomrule
\end{tabular}
\end{table}

\paperfig{\figpath/method_overview.png}
{One step of \method{}: the stage system couples the three constraint levels, the joint-coordinate
collocation rows and the Newton--Euler balance; the endpoint is reconstructed with the Gauss weights
and the exponential map and closed at velocity level.}
{fig:overview}

\subsection{Endpoint reconstruction and closure}

After the stage solve the endpoint is reconstructed with the Gauss weights,
\begin{align}
  r_{n+1} &= r_n+h\sum_{i=1}^3 b_i v_i, &
  v_{n+1} &= v_n+h\sum_{i=1}^3 b_i a_i, \label{eq:end-trans}\\
  Q_{n+1} &= Q_n\mathrm{Exp}\!\left(h\sum_{i=1}^3 b_i J_r^{-1}(\eta_i)\omega_i\right), &
  \omega_{n+1} &= \omega_n+h\sum_{i=1}^3 b_i \alpha_i .
  \label{eq:end-rot}
\end{align}
The reconstructed endpoint satisfies the position-level constraints only up to the quadrature
error ($O(h^7)$ per step, Lemma~\ref{lem:gauss-endpoint-defect}) and the velocity-level constraints
likewise. The endpoint closure map $\mathcal C_h$ then projects the endpoint velocities
$(v_{n+1},\omega_{n+1})$ onto the velocity-level constraint manifold of Eq.~\eqref{eq:fullva} at
$q_{n+1}$ by the minimal-norm correction, that is, by solving the KKT system with matrix
$\bigl[\begin{smallmatrix}I&\Phi_q^T\\ \Phi_q&0\end{smallmatrix}\bigr]$; positions and rotations are
not projected. The correction is an $O(h^7)$ perturbation of the step (Lemma~\ref{lem:endpoint-closure}).
The closed endpoint $x_{n+1}=\mathcal C_h(\mathcal E_h(Z))$, with $\mathcal E_h$ the map of
Eqs.~\eqref{eq:end-trans}--\eqref{eq:end-rot}, defines the one-step map
$\Psi_h(x_n)=x_{n+1}$.

\subsection{Newton solve and predictor}

The stage system $F_h(Z;x_n)=0$ is solved by Newton's method,
\begin{equation}
  D_ZF_h(Z^{(k)};x_n)\,\Delta Z=-F_h(Z^{(k)};x_n),\qquad Z^{(k+1)}=Z^{(k)}+\Delta Z,
  \label{eq:newton}
\end{equation}
with the $132\times132$ Jacobian assembled by forward-mode automatic differentiation of the single
residual function and factorized densely. The iteration stops when the residual norm or the update norm falls below the tolerance
$\tau_N=10^{-11}$ and is declared failed after $k_{\max}=80$ iterations. The predictor
extrapolates positions and velocities from the endpoint state with the constant endpoint
acceleration, sets the rotation increments, angular velocities and angular accelerations of the
stages to zero, and estimates the multipliers from the endpoint force balance; the consequences of
this choice for the convergence analysis are discussed after Lemma~\ref{lem:newton-envelope}. On
the chain the solve takes between five and six Newton iterations per step at every step size
(Section~\ref{sec:numerical}).

\begin{table}[htbp]
\centering
\caption{One step of \method{}.}
\label{tab:algorithm}
\small
\begin{tabular}{P{0.06\linewidth}P{0.84\linewidth}}
\toprule
Step & Operation \\
\midrule
1 & From $x_n=(q_n,v_n)$ and $h$, form the predictor $Z^{(0)}$. \\
2 & Assemble the stage residual $F_h(Z;x_n)$ of Eq.~\eqref{eq:g6fullva-expanded-stage}:
constraint rows at the three levels, joint-coordinate collocation rows, Newton--Euler rows. \\
3 & Iterate Eq.~\eqref{eq:newton} with the automatic-differentiation Jacobian until
$\|F_h\|\le\tau_N$ or $\|\Delta Z\|\le\tau_N$. \\
4 & Reconstruct the endpoint by Eqs.~\eqref{eq:end-trans}--\eqref{eq:end-rot}. \\
5 & Apply the velocity-level closure $\mathcal C_h$; record the Newton count and the endpoint
constraint norms. \\
\bottomrule
\end{tabular}
\end{table}

\paragraph{Cost}
One Newton iteration costs one residual and Jacobian evaluation and one dense $132\times132$
factorization; a step costs five to six such iterations. The Jacobian is block-sparse (each stage
couples to the others only through the collocation rows), which the present reference
implementation does not exploit. The two- and one-stage members of the same family, obtained by
replacing the Gauss tableau, have $88$ and $44$ unknowns and are used for the work/precision
comparison of Section~\ref{sec:numerical}.
